% ============================================================================
% Báo cáo BTL số 2 – Các vấn đề hiện đại KHMT, HK I 2026–2027
% Đề tài: Social Golfer Problem (SGP) bằng SAT, so sánh với CP/ILP
% Mẫu: Elsevier cas-sc (theo yêu cầu: \documentclass[a4paper,fleqn]{cas-sc})
% BIÊN DỊCH BẰNG XeLaTeX (Overleaf: Menu > Compiler > XeLaTeX), sau đó BibTeX.
% Mọi chỗ cần điền được đánh dấu \TODO{...} (chữ đỏ) – tìm "TODO" để rà soát.
% ============================================================================

% cas-sc nạp STIX kiểu Type 1 -> dễ sai ký hiệu toán với XeLaTeX; tắt gói này.
\makeatletter
\@namedef{ver@stix.sty}{9999/12/31 disabled for XeLaTeX}
\makeatother
\documentclass[a4paper,fleqn]{cas-sc}

\usepackage[vietnamese,english]{babel}
\usepackage{fontspec}
\setmainfont{texgyretermes}[
  Extension=.otf,
  UprightFont=*-regular,
  BoldFont=*-bold,
  ItalicFont=*-italic,
  BoldItalicFont=*-bolditalic]
\setsansfont{texgyreheros}[
  Extension=.otf,
  UprightFont=*-regular,
  BoldFont=*-bold,
  ItalicFont=*-italic,
  BoldItalicFont=*-bolditalic]
\setmonofont{texgyrecursor}[
  Extension=.otf,
  UprightFont=*-regular,
  BoldFont=*-bold,
  ItalicFont=*-italic,
  BoldItalicFont=*-bolditalic]
\usepackage{amsmath,amssymb}
\usepackage{booktabs}
\usepackage{graphicx}
\usepackage{xcolor}
\usepackage[numbers]{natbib}
\usepackage{xurl}

% cas-sc in footnote ORCID trống nếu không khai ORCID -> chỉ in khi có.
\ExplSyntaxOn
\let\OriginalPrintOrcid\printorcid
\RenewDocumentCommand\printorcid{}{
  \seq_if_empty:NF \g_stm_orcid_seq {\OriginalPrintOrcid}
}
\ExplSyntaxOff

\graphicspath{{figures/}}

% ---------- Thông tin nhóm (sửa tại đây) ----------
\newcommand{\TODO}[1]{\textcolor{red}{[TODO: #1]}}
\newcommand{\RepositoryURL}{https://github.com/thnhmai06/SAT-SGP}
\newcommand{\Affiliation}{Khoa Công nghệ thông tin, Trường Đại học Công nghệ, Đại học Quốc gia Hà Nội}
\newcommand{\EmailThanh}{\TODO{email Thành}}
\newcommand{\EmailCanh}{buidinhcanh2006@gmail.com}

\title[mode=title]{Giải bài toán Social Golfer bằng SAT và so sánh với CP, ILP}
\shortauthors{Mai Hải Thành, Bùi Đình Cảnh}
\shorttitle{SAT cho Social Golfer Problem}

\author[1]{Mai Hải Thành}
\cormark[1]
\ead{\EmailThanh}
\author[1]{Bùi Đình Cảnh}
\ead{\EmailCanh}
\address[1]{\Affiliation. Mã sinh viên: 24021627 (Mai Hải Thành), 24021392 (Bùi Đình Cảnh)}
\cortext[cor1]{Tác giả liên hệ}

\begin{document}

\begin{abstract}
\TODO{Tóm tắt 150--250 từ theo thứ tự: bài toán SGP; phương pháp SAT encoding và
cách giải (SAT nhiều lần / incremental nếu dùng); solver và cấu hình thí nghiệm
(CaDiCaL/Kissat/Glucose qua PySAT hoặc trực tiếp; so sánh Cplex CP, Cplex MIP,
Gurobi, CP-SAT); bộ instance $(g,s,w)$ và giới hạn thời gian; kết quả định lượng
chính (số instance SAT/UNSAT/timeout, so sánh với CP/ILP); kết luận ngắn.}
\end{abstract}

\begin{keywords}
Social Golfer Problem \sep SAT encoding \sep PySAT \sep Constraint Programming \sep
Integer Linear Programming \sep at-most-one
\end{keywords}

\maketitle

\section{Giới thiệu}
\TODO{Động cơ: SGP là bài toán tổ hợp kinh điển (lập lịch, thiết kế thí nghiệm, đối
xứng cao); vì sao SAT phù hợp; khoảng trống cần lấp. Nêu đóng góp theo danh sách
(i)--(iv): các encoding đã cài đặt, các solver đối chứng, quy trình thực nghiệm,
kết quả chính. Cuối mục: tóm tắt bố cục các mục sau.}

\section{Bối cảnh}
\subsection{Định nghĩa bài toán}
Có $n = g\cdot s$ người chơi golf. Trong mỗi tuần (round), họ được chia thành $g$
nhóm, mỗi nhóm đúng $s$ người, và mỗi người thuộc đúng một nhóm. Qua $w$ tuần,
không có hai người nào chơi cùng nhóm quá một lần. Bài toán thường ký hiệu
$g$-$s$-$w$; cần quyết định sự tồn tại của lịch $w$ tuần hợp lệ (bài toán quyết
định) và/hoặc tìm $w$ lớn nhất khi cố định $g,s$.

\TODO{Nêu rõ biến thể dùng trong BTL (quyết định hay tối ưu $w$) và quy ước ký
hiệu thống nhất với mã nguồn (Player, Group, Round, Match).}

\subsection{Cận trên và đối xứng}
\TODO{Cận trên $w \le \lfloor (n-1)/(s-1) \rfloor$; các đối xứng (hoán vị người
chơi, hoán vị nhóm trong tuần, hoán vị tuần) và cách phá đối xứng nếu có dùng.}

\subsection{Công trình liên quan}
\TODO{Các hướng CP, ILP, SAT cho SGP; bài báo thầy chỉ định về encoding
(\cite{csplib-sgp}); mã hóa cardinality \cite{sinz2005,prestwich2009}; PySAT
\cite{pysat-paper}.}

\section{Phương pháp SAT}
\subsection{Biến và ràng buộc}
\TODO{Định nghĩa biến, ví dụ $x_{p,r,k}$ = người $p$ thuộc nhóm $k$ ở tuần $r$.
Liệt kê và viết công thức cho: (C1) mỗi người đúng một nhóm mỗi tuần (ALO + AMO);
(C2) mỗi nhóm đúng $s$ người (cardinality); (C3) mỗi cặp gặp nhau tối đa một lần
(cần biến phụ cho ``cặp $(p,q)$ cùng nhóm ở tuần $r$''); (C4) phá đối xứng. Đưa
số biến, số mệnh đề theo $(g,s,w)$.}

\subsection{Các mã hóa AMO/ALO}
\TODO{Pairwise, Sequential Counter, Binary, Commander, Product (qua PySAT hoặc tự
cài). Bảng độ phức tạp kích thước (mẫu bên dưới).}

\begin{table}[htbp]
  \centering
  \caption{Kích thước mã hóa AMO cho nhóm $k$ literal \TODO{kiểm lại với cài đặt}.}
  \label{tab:amo}
  \begin{tabular}{lll}
    \toprule
    Mã hóa & Biến phụ & Mệnh đề \\
    \midrule
    Pairwise   & $0$ & $\binom{k}{2}$ \\
    Sequential & $\Theta(k)$ & $\Theta(k)$ \\
    Binary     & \TODO{} & \TODO{} \\
    Commander  & \TODO{} & \TODO{} \\
    Product    & \TODO{} & \TODO{} \\
    \bottomrule
  \end{tabular}
\end{table}

\subsection{Quy trình giải}
\TODO{Pipeline: sinh instance $\rightarrow$ encode CNF $\rightarrow$ gọi solver
$\rightarrow$ giải mã nghiệm $\rightarrow$ kiểm tra bằng checker độc lập. Nếu dùng
SAT nhiều lần / Incremental SAT khi tăng $w$ thì mô tả thuật toán (pseudo-code).}

\subsection{Tính đúng đắn}
\TODO{Lập luận ngắn: mô hình thỏa $\Leftrightarrow$ lịch hợp lệ; checker xác minh
mọi nghiệm SAT.}

\section{Mô hình đối chứng (CP và ILP)}
\TODO{Mô tả mô hình Cplex CP, Cplex MIP, Gurobi, CP-SAT (biến, ràng buộc, tham số
solver). Nhấn mạnh các mô hình biểu diễn cùng tập lịch hợp lệ.}

\section{Thiết kế thực nghiệm}
\subsection{Bộ dữ liệu}
\TODO{Danh sách instance $(g,s,w)$: CSPLib và các instance lớn (theo yêu cầu
``ma trận kích thước lớn''); lý do chọn.}

\subsection{Môi trường và giao thức}
\TODO{CPU/RAM/OS; phiên bản solver (CaDiCaL/Kissat/Glucose, Cplex, Gurobi,
OR-Tools); giới hạn thời gian mỗi ca; số lần lặp; số luồng; cách tính thời gian
(gồm cả thời gian encode); phân loại kết quả SAT/UNSAT/TIMEOUT/ERROR.}

\section{Kết quả và thảo luận}
% Hình vẽ bằng thư viện cnsplots (https://cnsplots.farid.one/latest/),
% xuất vào figures/ rồi \includegraphics.
\subsection{So sánh các SAT encoding}
\TODO{Bảng số biến/mệnh đề; biểu đồ thời gian giải; số instance giải được.}

\subsection{SAT so với CP và ILP}
\TODO{Bảng/biểu đồ tổng hợp SAT vs Cplex CP, Cplex MIP, Gurobi, CP-SAT; thảo luận.}

% Mẫu chèn hình:
% \begin{figure}[htbp]
%   \centering
%   \includegraphics[width=\linewidth]{ten_hinh.pdf}
%   \caption{...}\label{fig:ten}
% \end{figure}

\section{Hạn chế}
\TODO{Phạm vi instance, một lần chạy/ca, khác biệt cấu hình solver, license
academic, phần cứng.}

\section{Kết luận}
\TODO{Tóm tắt kết quả chính, hạn chế, hướng phát triển.}

\section*{Data availability}
Mã nguồn và kết quả thực nghiệm được lưu tại:
\url{\RepositoryURL}.

\section*{CRediT authorship contribution statement}
\textbf{Mai Hải Thành}: \TODO{vai trò, ví dụ Software, Methodology, Investigation}.
\textbf{Bùi Đình Cảnh}: \TODO{vai trò, ví dụ Writing -- original draft, Visualization, Validation}.

\section*{Declaration of competing interest}
Các tác giả tuyên bố không có xung đột lợi ích tài chính hay quan hệ cá nhân nào
có thể ảnh hưởng đến công trình được báo cáo.

\section*{Tuyên bố về công cụ AI hỗ trợ}
\TODO{Khai báo trung thực công cụ AI đã dùng (nếu có), mục đích và việc nhóm đã
rà soát nội dung.}

\bibliographystyle{cas-model2-names}
\bibliography{references}

\end{document}
